\section{Resultados de parámetros en formato wide}

\subsection{Parámetros estimados}

Las Tablas~\ref{tab:parametros_modelo_continuo} y~\ref{tab:parametros_modelo_dummy} presentan los
parámetros de las dos especificaciones candidatas. En ambas, \texttt{BIP} se mantiene como
alternativa de referencia para las variables comunes al conjunto de alternativas. Por lo tanto, para
variables como año, franja horaria, demanda, oferta, Censo, OSM y EOD, el coeficiente de
\texttt{BIP} se fija en cero y los coeficientes de \texttt{QR\_OTHER} y \texttt{QR\_RED} se
interpretan respecto de esa alternativa base.

Esta normalización no implica que dichas variables no estén asociadas a la elección de
\texttt{BIP}. Más bien, implica que el efecto se expresa en términos relativos: un coeficiente
positivo para \texttt{QR\_RED}, por ejemplo, indica que un aumento en la variable incrementa la
utilidad de \texttt{QR\_RED} respecto de \texttt{BIP}, manteniendo todo lo demás constante. Para las
variables que sí varían entre alternativas dentro de un mismo viaje, como tiempos y transbordos,
también se estima directamente un coeficiente para \texttt{BIP}.

Se presentan los coeficientes estimados y sus estadísticos \emph{t-ratio} robustos. Una versión
extendida de las tablas, que incluye errores estándar robustos, se deja en el
Anexo~\ref{sec:anexo_parametros_completos}. Los ceros reportados para \texttt{BIP} en variables
comunes corresponden a restricciones de normalización y no a parámetros estimados.

Además, para facilitar la interpretación de magnitudes, el Anexo~\ref{sec:anexo_odds_ratios}
presenta los \emph{odds ratios} asociados a ambas especificaciones. Estos se calculan como
\(\exp(\beta)\) o, en el caso de variables de tiempo medidas en segundos, como \(\exp(60\beta)\)
para reportar el cambio asociado a un minuto adicional.

\subsubsection{Modelo con proxy EOD continuo}

La Tabla~\ref{tab:parametros_modelo_continuo} presenta la especificación que incorpora la composición
histórica de ingreso EOD mediante proporciones continuas de hogares en los tramos
\texttt{ABC1\_proxy} y \texttt{D+E\_proxy}.

\begin{landscape}
\small
\setlength{\tabcolsep}{3pt}
\begin{longtable}{p{0.27\linewidth}rrrrrr}
\caption{Parámetros del modelo con proxy EOD continuo}
\label{tab:parametros_modelo_continuo}\\
\toprule
Variable &
\texttt{BIP} Est. & \texttt{BIP} t &
\texttt{QR\_OTHER} Est. & \texttt{QR\_OTHER} t &
\texttt{QR\_RED} Est. & \texttt{QR\_RED} t \\
\midrule
\endfirsthead

\toprule
Variable &
\texttt{BIP} Est. & \texttt{BIP} t &
\texttt{QR\_OTHER} Est. & \texttt{QR\_OTHER} t &
\texttt{QR\_RED} Est. & \texttt{QR\_RED} t \\
\midrule
\endhead

\bottomrule
\endfoot

\texttt{ASC} & -- & -- & -1.9172 & -23.32 & -3.5839 & -20.81 \\
\texttt{ANIO\_2025} & 0 & -- & 0.2965 & 33.32 & 0.2332 & 12.72 \\
\texttt{LAB\_PM} & 0 & -- & 0.0976 & 6.76 & 0.4604 & 15.25 \\
\texttt{LAB\_PT} & 0 & -- & 0.1571 & 8.87 & 0.2396 & 6.89 \\
\texttt{NO\_LAB} & 0 & -- & 0.1624 & 9.23 & 0.0510 & 1.35 \\
\texttt{LOG\_DEMAND} & 0 & -- & -0.0344 & -4.31 & -0.0671 & -4.24 \\
\texttt{LOG\_BUS\_STOP\_DENSITY} & 0 & -- & 0.0008 & 0.06 & -0.0910 & -2.92 \\
\texttt{LOG\_BUS\_LINE\_COUNT} & 0 & -- & 0.0164 & 1.13 & 0.0901 & 2.91 \\
\texttt{LOG\_METRO\_LINE\_COUNT} & 0 & -- & 0.0488 & 2.39 & 0.1235 & 3.14 \\
\texttt{T\_VEH} & 0.000020 & 0.66 & 0.000071 & 2.39 & -0.000006 & -0.19 \\
\texttt{T\_ESPERA\_INI} & -0.0017 & -29.21 & -0.0017 & -29.39 & -0.0014 & -20.08 \\
\texttt{T\_ESPERA\_TRASB} & -0.0003 & -3.59 & -0.0004 & -5.60 & 0.000045 & 0.60 \\
\texttt{N\_TRASB} & -0.0597 & -1.50 & -0.1248 & -3.15 & -0.1556 & -3.83 \\
\texttt{Educación univ. 18+} & 0 & -- & 0.0815 & 8.93 & 0.3073 & 16.47 \\
\texttt{Hacinamiento} & 0 & -- & -0.0652 & -5.66 & -0.0492 & -1.79 \\
\texttt{Discapacidad} & 0 & -- & 0.0028 & 0.24 & -0.0988 & -3.67 \\
\texttt{Inmigrantes} & 0 & -- & 0.1007 & 9.91 & 0.0865 & 3.68 \\
\texttt{Mujeres} & 0 & -- & -0.0002 & -0.06 & 0.0212 & 2.37 \\
\texttt{Asistencia parvularia} & 0 & -- & 0.0053 & 0.76 & 0.0487 & 3.11 \\
\texttt{Centros deportivos OSM} & 0 & -- & -0.0092 & -2.10 & 0.0078 & 0.89 \\
\texttt{Convenience OSM} & 0 & -- & 0.0135 & 3.07 & 0.0414 & 4.69 \\
\texttt{Playgrounds OSM} & 0 & -- & 0.0283 & 5.65 & 0.0662 & 6.44 \\
\texttt{Escuelas OSM} & 0 & -- & -0.0001 & -0.03 & -0.0284 & -2.54 \\
\texttt{Universidades OSM} & 0 & -- & -0.0197 & -9.15 & -0.0244 & -5.35 \\
\texttt{Refugios transporte OSM} & 0 & -- & -0.0306 & -4.65 & 0.0375 & 2.62 \\
\texttt{Accesos Metro OSM} & 0 & -- & -0.0096 & -2.83 & -0.0370 & -5.27 \\
\texttt{ABC1\_proxy EOD} & 0 & -- & -0.0154 & -1.87 & -0.0053 & -0.34 \\
\texttt{D+E\_proxy EOD} & 0 & -- & 0.0369 & 4.53 & -0.0123 & -0.67 \\

\end{longtable}
\end{landscape}

\subsubsection{Modelo con proxy EOD dummy}

La Tabla~\ref{tab:parametros_modelo_dummy} presenta la especificación que incorpora la composición
histórica de ingreso EOD mediante indicadores de alta concentración de hogares en los tramos
\texttt{ABC1\_proxy} y \texttt{D+E\_proxy}.

\begin{landscape}
\small
\setlength{\tabcolsep}{3pt}
\begin{longtable}{p{0.27\linewidth}rrrrrr}
\caption{Parámetros del modelo con proxy EOD dummy}
\label{tab:parametros_modelo_dummy}\\
\toprule
Variable &
\texttt{BIP} Est. & \texttt{BIP} t &
\texttt{QR\_OTHER} Est. & \texttt{QR\_OTHER} t &
\texttt{QR\_RED} Est. & \texttt{QR\_RED} t \\
\midrule
\endfirsthead

\toprule
Variable &
\texttt{BIP} Est. & \texttt{BIP} t &
\texttt{QR\_OTHER} Est. & \texttt{QR\_OTHER} t &
\texttt{QR\_RED} Est. & \texttt{QR\_RED} t \\
\midrule
\endhead

\bottomrule
\endfoot

\texttt{ASC} & -- & -- & -1.9081 & -23.16 & -3.5086 & -20.47 \\
\texttt{ANIO\_2025} & 0 & -- & 0.2963 & 33.30 & 0.2329 & 12.70 \\
\texttt{LAB\_PM} & 0 & -- & 0.0955 & 6.60 & 0.4529 & 14.97 \\
\texttt{LAB\_PT} & 0 & -- & 0.1522 & 8.60 & 0.2305 & 6.62 \\
\texttt{NO\_LAB} & 0 & -- & 0.1603 & 9.11 & 0.0447 & 1.18 \\
\texttt{LOG\_DEMAND} & 0 & -- & -0.0370 & -4.65 & -0.0722 & -4.56 \\
\texttt{LOG\_BUS\_STOP\_DENSITY} & 0 & -- & -0.0022 & -0.16 & -0.1134 & -3.60 \\
\texttt{LOG\_BUS\_LINE\_COUNT} & 0 & -- & 0.0251 & 1.72 & 0.1023 & 3.29 \\
\texttt{LOG\_METRO\_LINE\_COUNT} & 0 & -- & 0.0571 & 2.80 & 0.1251 & 3.19 \\
\texttt{T\_VEH} & 0.000020 & 0.67 & 0.000069 & 2.32 & -0.000008 & -0.25 \\
\texttt{T\_ESPERA\_INI} & -0.0017 & -29.34 & -0.0017 & -29.43 & -0.0014 & -20.04 \\
\texttt{T\_ESPERA\_TRASB} & -0.0003 & -3.58 & -0.0004 & -5.58 & 0.000046 & 0.61 \\
\texttt{N\_TRASB} & -0.0644 & -1.62 & -0.1285 & -3.24 & -0.1572 & -3.86 \\
\texttt{Educación univ. 18+} & 0 & -- & 0.0784 & 8.63 & 0.3194 & 17.16 \\
\texttt{Hacinamiento} & 0 & -- & -0.0556 & -4.84 & -0.0520 & -1.88 \\
\texttt{Discapacidad} & 0 & -- & 0.0123 & 1.09 & -0.1057 & -4.02 \\
\texttt{Inmigrantes} & 0 & -- & 0.1033 & 10.24 & 0.0962 & 4.09 \\
\texttt{Mujeres} & 0 & -- & 0.0005 & 0.15 & 0.0241 & 2.68 \\
\texttt{Asistencia parvularia} & 0 & -- & 0.0094 & 1.33 & 0.0672 & 4.17 \\
\texttt{Centros deportivos OSM} & 0 & -- & -0.0086 & -1.94 & 0.0113 & 1.29 \\
\texttt{Convenience OSM} & 0 & -- & 0.0119 & 2.74 & 0.0429 & 4.94 \\
\texttt{Playgrounds OSM} & 0 & -- & 0.0283 & 5.63 & 0.0644 & 6.26 \\
\texttt{Escuelas OSM} & 0 & -- & 0.0020 & 0.39 & -0.0262 & -2.35 \\
\texttt{Universidades OSM} & 0 & -- & -0.0178 & -8.03 & -0.0211 & -4.47 \\
\texttt{Refugios transporte OSM} & 0 & -- & -0.0304 & -4.58 & 0.0328 & 2.28 \\
\texttt{Accesos Metro OSM} & 0 & -- & -0.0104 & -3.06 & -0.0340 & -4.90 \\
\texttt{Alta concentración ABC1\_proxy} & 0 & -- & -0.0264 & -4.38 & -0.0421 & -3.44 \\
\texttt{Alta concentración D+E\_proxy} & 0 & -- & -0.0026 & -0.37 & -0.0110 & -0.65 \\

\end{longtable}
\end{landscape}

\subsection{Comparación general entre especificaciones}

Los parámetros de las variables base, temporales, de viaje, Censo y OSM son estables entre ambas
especificaciones. En particular, se mantienen los signos principales asociados a año, franjas
horarias, demanda local, tiempos de espera, educación universitaria, hacinamiento, inmigración,
refugios de transporte y accesos a Metro. Esto sugiere que la forma de representar el proxy EOD no
altera de manera sustantiva la estructura central del modelo.

La diferencia principal aparece en la lectura del bloque EOD. En la especificación continua,
\texttt{D+E\_proxy} presenta una asociación positiva y significativa con \texttt{QR\_OTHER}, mientras
que \texttt{ABC1\_proxy} muestra una señal negativa solo marginal. En la especificación dummy, la
alta concentración \texttt{ABC1\_proxy} presenta coeficientes negativos y significativos tanto para
\texttt{QR\_OTHER} como para \texttt{QR\_RED}, mientras que la alta concentración
\texttt{D+E\_proxy} no presenta una señal estadísticamente robusta.

La interpretación sustantiva de estos patrones se desarrolla en la sección siguiente. En esta sección
solo se deja establecido que ambas especificaciones son comparables, convergen correctamente y
conservan patrones centrales similares, pero capturan de forma distinta la dimensión histórica de
ingreso territorial proveniente de la EOD 2012.
